\documentclass[11pt]{article}
\usepackage[a4paper,margin=25mm]{geometry}
\usepackage[T1]{fontenc}
\usepackage{amsmath,amssymb,amsthm,mathtools,booktabs,longtable,hyperref}
\hypersetup{hidelinks}
\newtheorem{theorem}{Theorem}
\newtheorem{lemma}[theorem]{Lemma}
\newtheorem{proposition}[theorem]{Proposition}
\newtheorem{corollary}[theorem]{Corollary}
\newtheorem{remark}[theorem]{Remark}
\newcommand{\TV}{\operatorname{TV}}
\newcommand{\E}{\mathbb E}
\newcommand{\Prb}{\mathbb P}
\newcommand{\one}{\mathbf1}
\newcommand{\calC}{\mathcal C}
\newcommand{\calL}{\mathcal L}
\title{Global kinetic design for rare completion\\\large A direct scalar experimental oracle, a continuum optimizer, and its physical boundary}
\author{Independent Astra construction}
\date{5 November 2026}
\begin{document}
\maketitle
\begin{abstract}
We construct a global optimization operation for a rare desired first event.
The design variables are certified reciprocal kinetic multipliers which
preserve one unknown invariant distribution. The state space may be
exponential in the number of local variables. The construction chooses a
continuous allocation of the multipliers whose mean completion time from
equilibrium conditioned on not yet completing is within a prescribed
multiplicative factor of the global optimum. It uses local operating
experiments, a priced conditional-target sampler, gap certificates, and
ordinary linear optimization over the control constraints. It does not
acquire a full trajectory law or simulate an entire rare waiting time.
A direct short reverse no-return statistic estimates the normalized mean;
the common stationary target probability cancels from the optimization.
The normalized objective has a rarity-independent curvature bound, so
finite inexact Frank--Wolfe optimization needs polynomially many scalar
queries, rather than a grid exponential in the number of controls.
The kinetic multiplier interface is an explicit extra intervention premise;
it is not inferred from arbitrary field-dependent clocks. This is a
theoretical design capability, not a completed economical physical
catalyst. Bounded multipliers can improve mean completion time by at most
their largest allowed gain, so the actual rare completion exposure remains.
\end{abstract}

\section{The selected independent question}
Can local experiments select the globally best continuous allocation of
reversible kinetic acceleration, for a specified rare completion event,
without learning every operating rate, reconstructing the entire global
Markov chain, or generating complete first-event records?

The direction of the objective is essential. This document minimizes mean
time to a \emph{desired} event. It does not turn the maximization problem
of preventing a catastrophe into a convex problem. A finite menu of
preventive fields is a different previously earned operation. Nor is
minimum mean completion the same objective as minimum Kemeny constant,
fastest mixing, or minimum stationary target probability.

The result combines a new explicit interface with established machinery:
matrix-fractional convexity and reversible potential theory; the preserved
short no-return spectral statistic; finite experimental static and
operating updates; and inexact conditional-gradient optimization.
Historical priority for that combination is not established.

\section{Mathematical and experimental access}
\label{sec:access}
Let $\Omega$ be a finite, succinctly represented state space and let
$B\ne\varnothing$ be the known target. Write $A=\Omega\setminus B$,
$p=\pi(B)$ and $a=1-p$. Assume $0<p\le1/2$.
There are $d$ local generators $Q_j$. Each is reversible with respect to
the same full-support invariant law $\pi$. The admissible design family is
\begin{equation}
 Q(c)=\sum_{j=1}^d c_jQ_j,\qquad c\in\calC\subseteq[1,C]^d,\quad C\ge1.
 \label{eq:family}
\end{equation}
$\calC$ is a nonempty compact rational polytope with an explicitly priced
linear minimization operation. A budget polytope such as
$\sum_j w_j(c_j-1)\le b$ with box constraints is an example.
The baseline $\sum_jQ_j$ is irreducible. The number of channels, their
local supports, and the control constraints are input size; the number
of global states need not be polynomial.

Choose one rational
\[
 R\ge 2(C+1)\sum_j r_j,\qquad
 r_j\ge\max_x[-Q_j(x,x)],\qquad P(c)=I+Q(c)/R.
\]
The slightly enlarged numerical domain $[1,C+1]^d$ is only used for
one-sided finite differences. It requires no extra physical deployment
amplitude: its hypothetical updates can be composed from baseline
experimental channels as described below. Every $P(c)$ on this enlarged
domain is lazy and positive semidefinite. A certified normalized full-chain
gap $g\in(0,1]$ holds uniformly there. The baseline gap suffices since
each added Dirichlet form is positive semidefinite.

The operation is conditional on the following earned or supplied services.
\begin{enumerate}
\item A finite sampler for $\pi(\cdot\mid B)$ and, for the nonrare branch,
for $\pi(\cdot\mid A)$, with specified conditional total-variation error
and work. A relative count of $p$ is needed only when a supplied upper
bound does not establish the rare branch. The accepted correlated Gibbs
static construction supplies these at its original contrast, pinned
mixing, and local-degree scope. A singleton $B$ needs no conditional
target sampling; a certified simple upper bound on $p$ can avoid counting.
\item A requested-context baseline operating update of each $Q_j$, obtained
from faithful local copies with priced state pins, context clamps, actual
operating fields, rate bounds, detector and timer guarantees. The
preserved experiment-fed operating-coin construction supplies such an
update with a uniform conditional error. It does not transport a clock
from a balancing field to the operating field.
\item Structural closure, local Markov behavior, the common reversible
invariant, known supports, and numerical mixing/static margins. They are
premises. Local observations do not prove absence of hidden channels.
\item For \emph{deployment}, an actual reciprocal multiplier service
implementing \eqref{eq:family}, or a declared uniform multiplicative
conductance envelope around it. Arbitrary field changes, whose unknown
attempt clocks can depend on the field, do not supply this service.
Its response law, calibration, precision, slew, work and setup costs must
be supplied. Without it the conclusion is a design theorem for an
accessible generator family, not an intervention on the original object.
\end{enumerate}

There is an explicit way to generate hypothetical updates for any requested
$c\in[1,C+1]^d$ without relearning kinetics. Choose $j$ with weight
$r_j/\sum_i r_i$, retain the choice with probability $c_j/(C+1)$, and,
if retained, apply $I+Q_j/(2r_j)$. The resulting kernel is $P(c)$ when
$R=2(C+1)\sum_jr_j$; larger $R$ adds a known null probability.
Finite rational selectors and multiplier coins use an assigned error
budget. The local operating factory implements the requested baseline
channel. This is algebra on the observed baseline generator, not an
assumption about an unmeasured bias clock.

If the static generator is an unknown pairwise Gibbs model, balanced
local experiments supply the accepted static-weight and pinned-sampling
services. Their field metrology and repeated-cell conditional-error
premises remain unchanged. In particular, a sparse energy graph does not
prove sparse arbitrary kinetic context tables; this route queries contexts
as needed instead of enumerating those tables.

\section{The normalized objective and its geometry}
Let $\tau_B$ denote the physical first hitting time and define
\[
 M(c)=\E_{\pi(\cdot\mid A),Q(c)}\tau_B,
 \qquad f(c)=pR M(c).
\]
The factor $pR$ is common to every design. Minimizing $f$ and $M$ is
therefore exactly the same decision; no estimate of their astronomical
absolute scale is needed to interpret an optimizer.

Let $K(c)=-Q_A(c)/R$, where the subscript denotes killing on $B$,
and $v=\one_A$, in $\ell^2(A,\pi)$. Put $K_j=-Q_{j,A}/R$.
Then
\begin{equation}
 f(c)=\frac pa\langle v,K(c)^{-1}v\rangle_\pi,
 \qquad K(c)=\sum_jc_jK_j,\qquad K_j\succeq0.
 \label{eq:objective}
\end{equation}
All derivatives below are mathematical identities; the algorithm does
not form these exponentially large matrices.

\begin{theorem}[Rarity-free objective geometry]
For every $c\in[1,C+1]^d$,
\begin{align}
 \frac12&\le f(c)\le\frac1g,\label{eq:bounds}\\
 \partial_j f(c)&\le0,\qquad
 \sum_j c_j[-\partial_j f(c)]=f(c),\label{eq:homogeneous}\\
 0\le D^2f(c)[h,h]&\le\frac{2}{g}\|h\|_\infty^2.
 \label{eq:curvature}
\end{align}
In particular $f$ is convex, $\|\nabla f(c)\|_1\le1/g$, and its
Frank--Wolfe curvature on $\calC$ is at most
$2(C-1)^2/g$.
\end{theorem}
\begin{proof}
Zero extension of safe functions and Poincare give
$K(c)\succeq gpI$. Hence
$\langle v,K^{-1}v\rangle\le a/(gp)$, proving the upper bound.
Let $j_c=\langle v,K(c)v\rangle$ be stationary boundary flux in
normalized units. Cauchy--Schwarz gives
$\langle v,K^{-1}v\rangle\ge a^2/j_c$. Reversibility gives
$j_c\le p$ (in fact laziness gives $j_c\le p/2$), so $f\ge a\ge1/2$.

For $u=K^{-1}v$, differentiation gives
\[
 \partial_j f=-\frac pa\langle u,K_ju\rangle,
 \qquad
 D^2f[h,h]=\frac{2p}{a}\langle K(h)u,K^{-1}K(h)u\rangle,
 \quad K(h)=\sum_jh_jK_j.
\]
These prove convexity, negativity of coordinate derivatives, and
\eqref{eq:homogeneous}. Since $c_j\ge1$,
\[
 -\|h\|_\infty K(c)\preceq K(h)\preceq
 \|h\|_\infty K(c).
\]
Thus the self-adjoint operator $K^{-1/2}K(h)K^{-1/2}$ has norm at most
$\|h\|_\infty$. Substitution into the second derivative gives
$D^2f[h,h]\le2\|h\|_\infty^2 f(c)\le2\|h\|_\infty^2/g$.
For a chord $h=s-c$ in the control box, $\|h\|_\infty\le C-1$.
Taylor's integral remainder gives the stated curvature.
\end{proof}

The $\ell^\infty$ geometry in \eqref{eq:curvature} matters: the curvature
does not contain a sum over the number of global states, or an inverse
$p$. Control dimension still enters the cost of estimating a gradient.

\section{A direct scalar oracle with fixed update counts}
\label{sec:oracle}
Here the oracle returns $\widetilde f(c)$ satisfying
$|\widetilde f(c)-f(c)|\le\nu$ with conditional confidence $1-\delta_v$,
where $0<\nu\le1$. It returns no event time, terminal state or precursor
record. It does not construct the entire CDF. No global continuous
event clock is required: the mean of the continuous uniformized hitting
time equals the mean number of potential updates divided by $R$.

\subsection{One common branch for every control}
Set
\begin{equation}
 \xi=\frac{\nu g}{64},\quad
 L=\left\lceil\log_2\frac4{g\xi}\right\rceil,
 k=\lceil L/g\rceil,\quad p_{\rm gate}=\frac{\xi}{2k}.
 \label{eq:oracleSchedule}
\end{equation}
All ceilings are implementable by rational comparisons with powers of
two. Use a certified upper bound $p_U\le p_{\rm gate}$ if available.
Otherwise obtain a relative static interval $[p_L,p_U]$ with ratio at
most two. If $p_U\le p_{\rm gate}$ use the rare branch. Otherwise
$p_L>p_{\rm gate}/2$ and use the direct branch. The branch is shared by
all optimization queries because the invariant law is common.

\subsection{Rare branch: escape briefly from a deliberately supplied target}
Draw $Y\sim\pi(\cdot\mid B)$. Apply one $P(c)$ update. Return $Z=1$
exactly when this update leaves $B$ and the next $k$ updates stay in $A$.
Let $z(c)=\E Z$. This costs $k+1$ potential updates, even when the true
completion mean is exponentially large.

\begin{lemma}[Short scalar identity and approximation]
Under $p\le p_{\rm gate}$,
\[
 z(c)\ge g/4,
 \qquad \left|f(c)-\frac1{z(c)}\right|\le\nu/4.
 \label{eq:zApprox}
\]
\end{lemma}
\begin{proof}
The preserved killed spectral facts give eigenvalues
$0<\lambda=\lambda_1\le\lambda_2\le\cdots\le1$,
$\lambda\in[gp,p/a]$, $\lambda_2\ge g$, weights
$\alpha_i\ge0$ summing to one, and
$r=1-\alpha_1\le p/(ag)$. Reversibility gives the exact identity
\[
 pz=a\sum_i\alpha_i\lambda_i(1-\lambda_i)^k,
 \qquad f=p\sum_i\frac{\alpha_i}{\lambda_i}.
\]
Our gate implies $k\lambda\le\xi$ and $r\le\xi$.
The slow term divided by $a$ is at least
$(1-\xi)^2\lambda$. The sum of fast terms, divided by $a\lambda$,
is at most $e^{-gk}/(ag)\le\xi/2$. Consequently
\[
 (1-2\xi)\lambda\le pz/a\le(1+\xi)\lambda,
 \qquad 1-\xi\le (f/p)\lambda\le1.
\]
Thus $z\ge a(1-2\xi)g\ge g/4$ and
$(1-3\xi)a/z\le f\le(1+\xi)a/z$.
Since $a/z\le4/g$ and $p\le\xi$,
$|f-1/z|\le12\xi/g+4p/g\le16\xi/g=\nu/4$.
\end{proof}

Take $t=\nu g^2/128$. Make each implemented trial's conditional mean
differ from the ideal $z(c)$ by at most $t/2$, for example using a
conditional-target draw at TV error $t/4$ and each of $k+1$ updates at
error $t/[4(k+1)]$. Use
\[
 n_z=\left\lceil\frac{2}{t^2}\log\frac2{\delta_v}\right\rceil
\]
trials. Martingale Hoeffding, applied after conditioning on the previous
history, gives $|\bar Z-z|\le t$ with probability $1-\delta_v$.
No independence of repeated physical cells is asserted beyond their
declared conditional kernel guarantees. On this event $\bar Z\ge g/8$
and
\[
 \left|\frac1{\bar Z}-\frac1z\right|\le32t/g^2=\nu/4.
\]
Return an outward-rounded reciprocal, with rounding error at most
$\nu/4$. A zero or invalid empirical denominator returns a declared
failure value; it is outside the simultaneous good event. The result is
within $\nu$. The conservative statistical count is
$O(\nu^{-2}g^{-4}\log(1/\delta_v))$; it is not claimed optimal.

\subsection{Nonrare branch: a bounded direct mean}
Refine the single shared static count so
$|\widehat p/p-1|\le\rho=\nu g/64$, preserving a ratio-two interval.
Set
\[
 L_d=\left\lceil\log_2\frac{16}{\nu g}\right\rceil,
 J=\left\lceil\frac{L_d}{gp_L}\right\rceil,
 W=p_UJ\le\frac{2L_d}{g}+1.
\]
From $\pi_A$, let $H_B$ be the number of $P(c)$ updates to first hit $B$.
Estimate $\widehat p\E(H_B\wedge J)$ by bounded samples; replace $W$
by $2W$ if the chosen $\widehat p$ lies just outside the count interval.
The killed spectral bound gives
\[
 p\E(H_B-J)_+\le g^{-1}e^{-gpJ}\le\nu/16.
\]
The relative $p$ error contributes at most $\rho/g\le\nu/64$.
Choose each implemented trial's mean bias at most $\nu/16$, by giving
its entire finite trajectory TV error at most $\nu/(16\cdot2W)$;
allocate this to its initial draw and at most $J$ updates. Conditional
Hoeffding with $O(W^2\nu^{-2}\log(1/\delta_v))$ trials gives empirical
error at most $\nu/4$. Rational output rounding has its own $\nu/4$
allowance. This also supplies the desired oracle.

There is no hidden numerical $1/p$ cost in this branch:
$p_L>\xi/(4k)$ implies
$J\le4kL_d/(g\xi)+1$, a polynomial in $g^{-1},\nu^{-1}$ and logarithms.
Safe mass $a\ge1/2$ makes the finite safe-start rejection cost bounded.
All common static failures are charged once, separately from per-query
conditional confidence.

\section{Finite continuous design}
\label{sec:optimization}
Assume $0<\varepsilon\le1$ and put
\[
 \Delta=\varepsilon/2,\qquad W_c=\max(1,C-1),\qquad
 h=\frac{\Delta g}{32dW_c},\qquad
 \nu=\frac{\Delta h}{64dW_c}.
\]
Initialize at any rational $c_0\in\calC$. At iteration $t$ obtain
$d+1$ scalar values at $c_t$ and $c_t+he_j$, and set
\[
 G_{t,j}=\frac{\widetilde f(c_t+he_j)-\widetilde f(c_t)}h.
\]
Choose a feasible rational $s_t$ satisfying
\[
 G_t\cdot s_t\le\min_{s\in\calC}G_t\cdot s+\Delta/16,
\]
and use $c_{t+1}=(1-\gamma_t)c_t+\gamma_t s_t$ with
$\gamma_t=2/(t+2)$. An explicit rational linear program supplies this
operation in polynomial bit work; no exponentially sized list of designs
is generated.

\begin{theorem}[Global rare-completion allocation]
Take
\[
 T=\left\lceil\frac{16W_c^2}{g\Delta}\right\rceil.
\]
Give each of the $V=(d+1)T$ adaptive scalar calls conditional failure
probability at most $\delta/(2V)$ and give common acquisition, static
branching and margin certification total failure probability at most
$\delta/2$. On the simultaneous good event, exact rational control
arithmetic returns $c_T\in\calC$ satisfying
\[
 M(c_T)\le(1+\varepsilon)\min_{c\in\calC}M(c).
 \label{eq:global}
\]
All numerical control/output rounding can instead be assigned additive
objective allowance $\Delta/4$, with slack in the same conclusion.
The number of scalar calls is
$O(dW_c^2/(g\varepsilon))$; their required accuracy obeys
$\nu^{-1}=O(d^2W_c^2/(g\varepsilon^2))$.
\end{theorem}
\begin{proof}
Coordinate curvature gives
\[
 |G_{t,j}-\partial_jf(c_t)|\le h/g+2\nu/h.
\]
For any feasible competitor $s$, the inexact linear direction therefore
satisfies
\[
 \nabla f(c_t)\cdot(s_t-c_t)
 \le\nabla f(c_t)\cdot(s-c_t)+D,
\quad D\le dW_c(h/g+2\nu/h)+\Delta/16\le\Delta/8.
\]
Convexity and curvature $C_f\le2W_c^2/g$ imply, for
$e_t=f(c_t)-\min f$,
\[
 e_{t+1}\le(1-\gamma_t)e_t+\gamma_tD+
                 \tfrac12\gamma_t^2C_f.
\]
The usual elementary induction, starting with $\gamma_0=1$, gives
$e_T\le D+2C_f/(T+2)\le3\Delta/8$. This is less than $\Delta$ even
after the declared rounding allowance. Since $\min f\ge1/2$ and
$\Delta=\varepsilon/2$, the additive bound implies
$f(c_T)\le(1+\varepsilon)\min f$. The positive constant $pR$ cancels.
Conditional union bounds apply to the adaptive calls; their input points
may depend on every earlier result. The required uniform conditional
oracle guarantees were constructed above.
\end{proof}

\paragraph{A stopping certificate.}
At any iterate, let $s_t$ be the approximate linear minimizer. If
$\eta_G\ge\|G_t-\nabla f(c_t)\|_1$ is certified, then
\[
 0\le f(c_t)-\min f
 \le G_t\cdot(c_t-s_t)+(C-1)\eta_G+\Delta/16.
\]
The right side is an observable inexact Frank--Wolfe gap. It can stop an
instance earlier than the sufficient worst-case iteration schedule, but
its value is meaningful only with the same scalar-oracle confidence
and finite arithmetic allowances.

\paragraph{Bit and memory costs.}
The controls, linear programs and value outputs are rational. Exact
convex combinations have bit length polynomial in $T$ and the input
bit length. For finite arithmetic one can round to a feasible rational
point with $\ell^\infty$ displacement $\eta_c$; the objective changes by
at most $\eta_c/g$. A feasible rounding operation must be provided by
the actual control polytope, not assumed for an arbitrary nonlinear
constraint set. Working memory comprises a control vector, local state,
counters, and the inherited static-sampling data structure. No vector of
length $|\Omega|$ or long event record is needed.

\section{Execution costs and a real physical deployment statement}
\label{sec:physical}
Let $C_B(\theta)$, $C_A(\theta)$ and $C_P(\theta)$ denote the complete
costs of a target draw, safe draw and operating update at the requested
conditional error, respectively. These include all relevant physical
experiments; they are not unit-cost oracles. In the rare branch one scalar
call costs at most
\[
 n_z\left[C_B(t/4)+(k+1)C_P\!\left(\frac{t}{4(k+1)}\right)
               +C_{\rm predicate}+C_{\rm bookkeeping}\right].
\]
In the nonrare branch it costs
\[
 O\!\left(W^2\nu^{-2}\log(1/\delta_v)\right)
 \left[C_A(\theta_A)+J C_P(\theta_P)+J C_{\rm predicate}
                                      +C_{\rm bookkeeping}\right],
\]
where $\theta_A+J\theta_P\le\nu/(32W)$ is sufficient.
Multiply by $V$. Add common static, acquisition and margin checks,
linear optimization, output and deployment. This is an explicit
polynomial reduction, not an economical numerical bill.

For the inherited one-coordinate operating factory, $C_P(\theta)$
contains $O(\log(1/\theta))$ priced races, each with a pin and release,
$O(\rho^{-1}\log^2(1/\theta))$ capped dwell, and the declared timing
resolution of order $\theta/[\rho\log(1/\theta)]$, with its conservative
constants. Those costs may be large. This theorem changes what is
estimated and the dimension of the design search; it does not claim
that a physical race is cheap, or that its sufficient timer requirement
is a universal lower bound. There is no cost of simulating an entire
mean lifetime or constructing a full marked path. The model's actual
completion runtime is still $M(c_T)$ and must be paid when the desired
event is executed.

\subsection{Finite preparation and finite slew}
Suppose a full-system repairable implementation can run through target
states during preparation, a state-mass floor $\zeta>0$ and physical
baseline gap $\gamma>0$ are certified, and the state can be inspected at
a declared completion time. Baseline evolution for
\[
 t_{\rm mix}\ge\gamma^{-1}\log(1/(\eta\zeta))
\]
gives a pointwise density ratio in $[1-\eta,1+\eta]$ relative to $\pi$
from any start, using the standard finite-state $L^2$ to pointwise bound.
Allow surplus for finite timing and gap calibration. Ramp from baseline
to the chosen $c_T$ with the actual finite slew. Every intermediate
generator preserves $\pi$, so the same pointwise envelope is preserved
by the time-inhomogeneous transition kernel during the ramp. This
argument does not require a free instantaneous switch or a globally
frozen network.

Inspect and retain a state in $A$. Since $a\ge1/2$, one trial succeeds
with probability at least $(1-\eta)/2$. On success the initial law
$\nu_{c_T}$ obeys
\[
 e^{-\beta}\le\frac{d\nu_{c_T}}{d\pi_A}\le e^\beta,
 \qquad \beta=\log\frac{1+\eta}{1-\eta}.
\]
Finite retries with this success floor give explicit logarithmic caps
and preparation abort probability. The target must be nondestructive
until the declared mission starts; irreversible destruction during
preparation is outside this operation. Only the selected design needs
physical deployment/preparation. Acquisition and computation happen
before the mission unless a separately priced holding mode is supplied.

Because the mean is unbounded as a random loss, an abort cannot simply
be charged as an additive total-variation risk. The theorem reports the
conditional successful-start mean and the abort probability separately.
Alternatively, an uncapped repeat-until-success protocol has a bounded
expected preparation overhead, which must be added to the completion
mean. A finite hard cap produces an explicit abort outcome.

\subsection{Deployment uncertainty}
If the actual deployed generator $\widetilde Q(c)$ preserves $\pi$ and
its Dirichlet form is between $e^{-\kappa}$ and $e^\kappa$ times that
of $Q(c)$ uniformly in $c$, inverse Loewner order gives
\[
 e^{-\kappa}M(c)\le\widetilde M(c)\le e^\kappa M(c).
\]
If every successful preparation has the above $e^{\pm\beta}$ density
envelope, the actual design decision obeys
\[
 \E_{\nu_{c_T},\widetilde Q(c_T)}\tau_B
 \le e^{2(\kappa+\beta)}(1+\varepsilon)
   \inf_{c\in\calC}\E_{\nu_c,\widetilde Q(c)}\tau_B.
\]
This is a uniform robustness statement, not a finite-observation proof
of the multiplier response law. Detector, clock and field errors during
acquisition enter $C_P$'s conditional error; actuation calibration enters
$\kappa$. A directed error breaking detailed balance or changing the
invariant law needs a separately proved envelope. It cannot be counted
as the same $\kappa$ by assertion.

\section{What this makes possible, and what it cannot make possible}
\begin{corollary}[No exponential physical speedup from bounded reciprocal gains]
If $\one\in\calC$ and all feasible coordinates satisfy $1\le c_j\le C$,
then
\[
 \frac{M(\one)}C\le M(c)\le M(\one).
\]
\end{corollary}
\begin{proof}
$K(\one)\preceq K(c)\preceq C K(\one)$, and inverse order reverses
both inequalities in \eqref{eq:objective}.
\end{proof}
An exponentially rare equilibrium target can therefore remain
exponentially slow to realize. The capability earned here is global
design and experimental decision at a cost independent of that runtime,
not free physical realization of the rare product. Exponential gains
require exponentially large allowed kinetic gains, a changed invariant,
a different initialization, a new transition mechanism, or another
resource outside this family. This corollary is a direct consequence of
classical form comparison, not a claim of a new universal law of
catalysis.

The target family is meaningful beyond a tiny fixture. For $N$ binary
variables with bounded-degree interacting Gibbs weights, fixed
conditional-surprisal contrast and certified pinned uniqueness, the
accepted static construction supplies $B$ as all-terminal disconnection;
the present theorem then designs the fastest deliberate stress-test
completion within its reciprocal kinetic family. For any singleton
target configuration, $\pi_B$ is deterministic. If all one-site
conditional probabilities of its specified bits are at most
$q_*<1$, then $p\le q_*^N$. For sufficiently large $N$ this certifies
the rare branch without estimating $p$ at all. Correlated Gibbs
interaction and unknown baseline local clocks are allowed under the
declared certificates. In both examples $|\Omega|=2^N$ while the
control dimension and oracle reduction are polynomial in local input
size and numerical inverse margins.

\section{Strongest applicable comparisons and honest scope}
Reversible hitting-time variational principles are classical
\cite{aldous}. Convex design of explicit Markov chains, including
fastest mixing and effective-resistance objectives, is established
\cite{boyd}; reversible Kemeny optimization has an explicit semidefinite
formulation \cite{patel}. Inexact Frank--Wolfe and its gap certificate are
also established \cite{jaggi}. This work does not claim these ingredients
as inventions.

The completed difference is the \emph{succinct unknown-generator
interface}: short local experimental updates plus conditional target
draws yield a scalar objective oracle with finite confidence, then a
global continuum choice, without a matrix with one row per global state,
uniform kinetic table recovery, a control grid, or full marked-record
certification. An explicit-state SDP has at least exponential input
size on the $2^N$ family; ordinary direct physical hitting-time sampling
has exposure proportional to the mean lifetime. These are comparisons
of representations under the same supplied local access, not claims
that every specialist method is inferior. A method using the same
short-return reduction and convex geometry can match this capability.

The accepted whole-CDF machinery already removes inverse-rarity
computation in its model. This construction does not claim another
exponential rarity speedup over it; it eliminates unnecessary whole-CDF
and marked-law outputs and replaces finite-menu enumeration by a global
continuous optimization theorem. On independent product models,
factorization gives especially strong ordinary integral formulas; the
executed diagnostic uses those as an independent reference rather than
pretending a general sampler is the best solver for that subclass.

Arbitrary field-dependent clocks invalidate the affine generator family.
Finite experiments at finitely many fields cannot certify an unassumed
uniform response over a continuum. Maximizing the convex objective for
catastrophe prevention is a different optimization problem. Adaptive
feedback policies, all-working initial states, irreversible setup,
unknown hidden mechanisms, nonreversible driving, and changing invariant
laws are not supplied by this result.

\begin{proposition}[Why the response premise cannot be replaced by finite queries]
Consider a one-bit model at fields $u\in[0,1]$ with fully known logit $u$
and rates $r(u)\sigma(u)$ from zero and $r(u)[1-\sigma(u)]$ from one.
Let $r$ range over smooth functions in $[1,C]$, with no supplied modulus
of continuity or derivative bound. Even an oracle that returns the exact
$r(u)$ at each queried field cannot, with a uniformly bounded number of
queries, guarantee a global constant-factor approximation to minimum
mean completion time over this class when the desired factor is smaller
than $C\sigma(0)/\sigma(1)$. This is a finite-query statement, not a
lower bound for an unpriced infinitely resolving continuous sweep.
\end{proposition}
\begin{proof}
Place $m$ disjoint open intervals in $[0,1]$. In instance $j$, set
$r_j(u)=1+(C-1)b_j(u)$ where $b_j$ is a smooth bump supported in interval
$j$ and equal to one at its center. Until an algorithm queries inside
the selected interval, its transcript is identical to that for $r=1$.
For a uniform random $j$, a randomized algorithm making at most $q$
queries and outputting one field intersects the selected interval with
probability at most $(q+1)/m$: follow its transcript on the baseline,
then count the intervals containing one of its queries or its output.
If it misses, its completion rate is at most $\sigma(1)$ whereas the
optimum is at least $C\sigma(0)$. Hence its mean-time approximation
factor is at least $C\sigma(0)/\sigma(1)$. Letting $m$ grow defeats any
uniform finite query bound. Faithful finite experiments are no stronger
than this exact response oracle.
\end{proof}

The known affine multiplier law in \eqref{eq:family} is therefore a
substantive structural intervention resource. A finite calibration with
a proved response modulus can instead be charged at that modulus's
scale; the proposition does not forbid such a construction.

\section{Execution record}
The accompanying code checks the objective and derivative identities on
small independent chains by two representations (explicit killed linear
algebra and a product integral), and performs continuum allocation on a
growing independent family using the strongest available factorized
integral as an audit oracle. It records numerical optimization gaps and
target rarity. These are deterministic mathematical diagnostics under
known synthetic parameters. They do not execute the general correlated
static service, the conservative finite experimental schedule, a physical
kinetic multiplier, or a real rare completion. Details and any failed
checks are retained in the execution record. Numerical quadrature errors
are separately reported; those floating diagnostics are not a substitute
for the finite rational theorem.

\begin{center}
\begin{tabular}{rccc}
\toprule
$N$ & Target mass $p$ & Relative numerical FW gap & Mean after design\\
\midrule
8 & $8.78\cdot10^{-3}$ & $3.92\cdot10^{-5}$ & $18.02$\\
16 & $1.25\cdot10^{-4}$ & $1.19\cdot10^{-5}$ & $516.76$\\
32 & $7.77\cdot10^{-9}$ & $1.99\cdot10^{-5}$ & $3.99\cdot10^6$\\
64 & $1.22\cdot10^{-16}$ & $1.63\cdot10^{-5}$ & $1.21\cdot10^{14}$\\
128 & $5.33\cdot10^{-32}$ & $4.35\cdot10^{-6}$ & $1.36\cdot10^{29}$\\
\bottomrule
\end{tabular}
\end{center}
The state space at $N=128$ has $2^{128}$ states. The improvement over equal
feasible allocation ranges only from $1.096$ to $1.233$ across these
examples, as expected under bounded kinetic gains. An ordinary constrained
optimizer on the same product integral agrees within the numerical gap;
this is not a separation from that strong subclass solver. Six small
explicit killed-matrix checks agree with the product formula and its
derivative homogeneity identity. A separate interacting three-bit
90-decimal-digit check has $p=3.13\cdot10^{-16}$ and $k=1242$;
$f=4.022433879193464\ldots$ and $1/z=4.022433879193855\ldots$,
an error $3.90\cdot10^{-13}$ below the declared allowance $0.005$.

The sufficient general experimental schedule is \emph{not} the fast
product calculation. For the $N=128$ instance, $\varepsilon=0.1$ and
$\delta=0.05$, the elementary theorem schedule gives approximately
$4.12\cdot10^8$ scalar calls at $\nu=7.46\cdot10^{-15}$ and
$10^{61.65}$ potential operating updates before the per-update physical
factory bill. It is evaluated, not executed. The absence of an inverse
global-rarity term is therefore a theoretical scaling result, not an
earned useful experimental regime. Neither these constants nor this
schedule are lower bounds. A practical implementation would need a
different gradient/oracle realization or stronger instance structure;
that improvement is not claimed here.

\begin{thebibliography}{9}
\bibitem{accepted} Preserved Astra acquired first-event package,
\texttt{acquired\_first\_event\_research\_package}, including
\texttt{next5\_decision/experimental\_operating.tex},
\texttt{experimental\_static.tex}, and the preceding finite CDF proof.
Every file is retained unchanged; hashes are in the companion record.
\bibitem{aldous} D. Aldous and J. Fill, \emph{Reversible Markov Chains and
Random Walks on Graphs}, Chapter 3.7,
\url{https://www.stat.berkeley.edu/~aldous/RWG/Book_Ralph/Ch3.S7.html}.
\bibitem{boyd} S. Boyd, \emph{Convex Optimization of Graph Laplacian
Eigenvalues}, Proceedings of the International Congress of Mathematicians
3:1311--1319 (2006),
\url{https://stanford.edu/~boyd/papers/cvx_opt_graph_lapl_eigs.html}.
\bibitem{patel} R. Patel, P. Agharkar, and F. Bullo,
\emph{Robotic Surveillance and Markov Chains with Minimal Weighted Kemeny
Constant}, IEEE Transactions on Automatic Control 60(12):3156--3167
(2015), DOI 10.1109/TAC.2015.2426317. Author manuscript:
\url{https://motion.me.ucsb.edu/pdf/2014b-pab.pdf}.
\bibitem{jaggi} M. Jaggi, \emph{Revisiting Frank--Wolfe: Projection-Free
Sparse Convex Optimization}, ICML, PMLR 28(1):427--435 (2013),
\url{https://proceedings.mlr.press/v28/jaggi13.pdf}.
\end{thebibliography}
\end{document}
